\centerline{\bf \Huge Разнобой}

\begin{flushright}
        \scriptsize{Этот листочек составлен с помощью чата гпт}
\end{flushright}


\problem{1}
На асфальте нарисована полоса $1 \times 10$ для игры в «классики». Из центра первого квадрата надо сделать 9 прыжков по центрам квадратов (иногда вперёд, иногда назад) так, чтобы побывать в каждом квадрате по одному разу и закончить маршрут в последнем квадрате. Аня и Варя обе прошли полосу, и каждый очередной прыжок Ани был на то же расстояние, что и очередной прыжок Вари. Обязательно ли они пропрыгали квадраты в одном и том же порядке?

\problem{2}
Даны окружности $\omega_1$ и $\omega_2$. Пусть $M$ - середина отрезка, соединяющего их центры. На $\omega_1$ выбрана точка $X$, а на $\omega_2$ - точка $Y$ так, что $M X=M Y$. Найдите геометрическое место середин отрезков $X Y$.

\problem{3}
На прямой дороге стоят школа и дома Ани и Бори. Каждый день Аня выходит из дома в 8:00 и идет в школу. Однажды Боря выбежал из дома в школу в 8:00 и догнал Аню за 30 минут. На следующий день он выбежал в 8:10 и догнал Аню за 40 минут. В какое время ему надо выбежать, чтобы встретить Аню на выходе из её дома? (Скорость Ани всегда постоянна, скорость Бори тоже постоянна.)

\problem{4}
Решите уравнение $\sqrt{x+4}-\sqrt{6-x}+4=2 \sqrt{24+2 x-x^2}$.

\problem{5}
На часах три стрелки, каждая вращается в ту же сторону, что и обычно, с постоянной ненулевой, но, возможно, неправильной скоростью. Утром длинная и короткая стрелки совпали. Ровно через 3 часа совпали длинная и средняя стрелки. Еще ровно через 4 часа совпали короткая и средняя стрелки. Обязательно ли когда-нибудь совпадут все три стрелки?

\problem{6}
Квадратный трёхчлен $f(x)=a x^2+b x+c$ имеет два различных вещественных корня $x_1$ и $x_2$. Известно, что $f\left(x_1+x_2\right)=2025$. Чему может равняться $c$?

\problem{7}
Решите систему уравнений

$$
\left\{\begin{array}{l}
\frac{2}{x^2+y^2}+x^2 y^2=2, \\
x^4+y^4+3 x^2 y^2=5 .
\end{array}\right.
$$